% TOP_PROBLEM: {"id": 30006803, "title": "All-Pairs Shortest Paths in Truly Subcubic Time", "category_id": 15, "category": {"id": 15, "name": "computer_science", "display_name": "Computer Science", "description": "Computational complexity, algorithms, and theoretical CS.", "slug": "computer-science", "order_index": 15, "created_at": "2026-07-31T15:26:25.671Z"}, "status": "open"}
% FIRST_POSTED: 2026-09-27
% !TeX program = lualatex
% Compile twice: lualatex 30006803.tex
% All references and prose are embedded; no BibTeX database or external inputs.
\documentclass[11pt,a4paper]{article}
\usepackage[margin=24mm,headheight=14pt]{geometry}
\usepackage{fontspec}
\setmainfont{Latin Modern Roman}
\setsansfont{Latin Modern Sans}
\setmonofont{Latin Modern Mono}
\usepackage{amsmath,amssymb,mathtools}
\usepackage{unicode-math}
\setmathfont{Latin Modern Math}
\usepackage{microtype}
\usepackage{enumitem,xurl,needspace}
\usepackage[dvipsnames]{xcolor}
\usepackage{hyperref}
\hypersetup{unicode=true,colorlinks=true,linkcolor=MidnightBlue,urlcolor=MidnightBlue,
 pdftitle={Research Notebook: Section 30006803},pdfauthor={Research notebook},bookmarksnumbered=true}
\usepackage{bookmark}
\usepackage{tocloft}
\setlength{\cftsecnumwidth}{3em}
\cftsetpnumwidth{2.5em}
\cftsetrmarg{4em}
\usepackage{titlesec}
\titleformat{\section}{\Large\bfseries\raggedright}{\thesection}{0.7em}{}
\usepackage{fancyhdr}
\pagestyle{fancy}\fancyhf{}
\fancyhead[L]{\small\sffamily Open Problems | Research notebook}
\fancyhead[R]{\small 232 | 27 September 2026}
\fancyfoot[C]{\thepage}
\setlength{\parindent}{0pt}
\setlength{\parskip}{5pt plus 1pt}
\setlength{\emergencystretch}{3em}
\setcounter{tocdepth}{1}
\setcounter{secnumdepth}{1}
\setlist[itemize]{leftmargin=3em,itemsep=5pt,topsep=4pt,parsep=0pt}
\allowdisplaybreaks
\newcommand{\N}{\mathbb{N}}
\newcommand{\Z}{\mathbb{Z}}
\newcommand{\Q}{\mathbb{Q}}
\newcommand{\R}{\mathbb{R}}
\newcommand{\C}{\mathbb{C}}
\newcommand{\F}{\mathbb{F}}
\newcommand{\A}{\mathbb{A}}
\newcommand{\PP}{\mathbb P}
\newcommand{\E}{\mathbb{E}}
\newcommand{\eps}{\varepsilon}
\newcommand{\dd}{\,\mathrm d}
\newcommand{\sref}[2]{\hyperlink{source:#1:#2}{[S#2]}}
\newcommand{\eref}[2]{\hyperlink{extra:#1:#2}{[E#2]}}
% Turn the next switch off for a shorter reading copy. The quotations preserve
% every exactTarget field, including qualifications omitted from a short summary.
\newif\ifshowcatalogue
\showcataloguetrue
\newcommand{\cataloguescope}[1]{\ifshowcatalogue
 \paragraph{Exact scope in the supplied catalogue.}
 \begin{quote}\small #1\end{quote}\fi}

% Independently compilable extraction; original section text retained.
\numberwithin{equation}{section}
\begin{document}
\setcounter{section}{0}
% ================================================================
% problemId: 30006803
% Canonical problem ID: 30006803
% exactTarget SHA-256: 46ce2b3fa756a88aa7a3d24bc455a35bf0b9901441592a17ed35094d5854e1bd
\clearpage
\section*{\texorpdfstring{All-Pairs Shortest Paths in Truly Subcubic Time}{All-Pairs Shortest Paths in Truly Subcubic Time}}\label{30006803}
\subsection{Definitions and mathematical statement}
For a weighted graph $G=(V,E,w)$, the distance $d(u,v)$ is the infimum of total edge weight over paths from $u$ to $v$, with $+\infty$ for unreachable pairs. In the conventional finite-distance formulation there are no relevant negative-weight cycles. Exact all-pairs shortest paths outputs the whole distance matrix, not approximations. The requested running-time improvement is
\[
 \exists\eps>0\quad T(n)=O(n^{3-\eps})\quad(n=|V|),
\]
with a fixed exponent saving. \textbf{Model qualification:} the catalogue does not pin down directionality, weight range or encoding, and the computational model. Those choices are necessary for a unique complexity assertion and must be fixed from the intended source; unrestricted bit-length input cannot have a bound depending only on $n$. Approximate distances do not satisfy the stated exact target.
\par\smallskip\textit{Formulation and concept references:} \sref{30006803}{1}.
\cataloguescope{Determine whether exact all-pairs shortest paths on weighted n-vertex graphs can be computed in O(n\textasciicircum{}(3-epsilon)) time for some fixed epsilon greater than zero.}
\subsection{Short English statement}
Can exact distances between every pair of vertices be computed with a fixed polynomial improvement over cubic time, once the weighted-graph and machine conventions are fixed?
% BEGIN RESEARCH ATTEMPT 232
\subsection{Research attempt}

\paragraph{Current frontier and model audit.}
The full manuscript behind the catalogue citation treats approximation
algorithms and explicitly contrasts exact weighted APSP with the easier
unweighted and approximate settings. In particular, its Introduction
already regards weighted \emph{undirected} graphs as part of the exact
barrier. \sref{30006803}{1}; \eref{30006803}{1}, Abstract and Introduction.
Its stated matrix-multiplication exponent is a date-specific value, not
something this notebook treats as a current constant. The 2024--2026
search in this pass did not provide a verified exact algorithm resolving
the unrestricted weighted question.

\textbf{Editorial scope note.} The original model qualification remains
in force. Neither a claim for unrestricted bit strings in time depending
only on $n$, nor a unit-cost operation on arbitrarily long integers, is
silently selected as the intended target. The calculations below specify
their own finite-integer auxiliary problems and show exactly which
reductions they support. They do not repair the missing global convention
by decree.

\paragraph{Attack route.}
One route is to encode min-plus multiplication inside ordinary matrix
multiplication. It needs a representation whose size depends polynomially
on the \emph{bit length} of weights rather than their magnitude. A second
tries to round sufficiently accurate approximations to exact distances;
it needs an error scale compatible with the running-time guarantee. A
third seeks to exploit undirectedness. The attempt gives a positive-weight
undirected gadget that still encodes unrestricted finite min-plus data,
then quantifies the representation and rounding obstructions. This
identifies a concrete missing step rather than confusing a logarithmic
speedup, bounded-weight subcase, or approximation with the required fixed
exponent saving.

\paragraph{Attempt.}
\textbf{1. An unavoidable encoding check.}
In a bit-operation model, take a two-vertex graph whose only edge has a
positive integer weight of $B$ bits. Its nonzero exact distance has $B$
bits too, so writing the output alone takes $\Omega(B)$ bit operations.
Allowing arbitrary $B$ excludes a running-time bound depending only on
$n$. This observation is not a lower bound against the customary
word-RAM or real-RAM APSP conjectures; it is a reason their weight and
machine conventions must be specified. Outputting a symbolic pointer to
an input weight would itself change the output encoding.

For the next steps fix two $n\times n$ matrices $A,B$ with entries in
$\{0,1,\ldots,W\}$ and define their exact min-plus product by
\[
 (A\star B)_{ij}=\min_{1\leq k\leq n}(A_{ik}+B_{kj}).
\]
All subsequent integers have their ordinary binary encoding unless an
explicit unit-cost arithmetic interpretation is stated.

\textbf{2. Undirected positive weights still encode min-plus products.}
Create three disjoint vertex layers $I=\{i_1,\ldots,i_n\}$,
$K=\{k_1,\ldots,k_n\}$ and $J=\{j_1,\ldots,j_n\}$. Let
$H=2W+1$. Put undirected edges
\[
 w(i_i,k_k)=H+A_{ik},\qquad w(k_k,j_j)=H+B_{kj},
\]
and no others. Every two-edge $i_i$--$j_j$ path has weight
$2H+A_{ik}+B_{kj}\leq2H+2W$. A path between those layers has an even
number of edges. If it is not a two-edge path, it has at least four,
and hence weight at least $4H>2H+2W$. All weights are positive, so
removing cycles cannot increase path length. It follows exactly that
\begin{equation}\label{30006803:gadget}
 d(i_i,j_j)=2H+(A\star B)_{ij}.
\end{equation}
The graph has $3n$ vertices and $2n^2$ edges, and its weights use
$O(1+\log(W+1))$ bits. Subtracting the known $2H$ recovers the entire
min-plus product. Thus undirectedness does not eliminate the min-plus
subproblem, even when all graph weights are positive. No negative cycle,
approximate path, or implicit exponentially large graph is used.

Conversely, if a min-plus product algorithm applies uniformly to the
necessary weight range and to $+\infty$, take the minimum of the identity
matrix (zero diagonal, $+\infty$ elsewhere) and a graph's adjacency
matrix. Repeated min-plus squaring computes minimum weights of walks
with at most $2^r$ edges. In a graph without negative cycles every finite
shortest distance has a simple-path representative with at most $n-1$
edges. Therefore $O(\log n)$ squarings suffice for APSP. The same argument
applies in particular to positive-weight undirected graphs. For bounded
integer inputs, intermediate finite weights need only $O(\log n+\log
(W+1))$ bits when the relevant shortest paths are retained.
A time bound $O(n^{3-\eta})$ for those products would give
$O(n^{3-\eta}\log n)=O(n^{3-\eta/2})$ for fixed $\eta>0$ and
sufficiently large $n$, after accounting for the same model's arithmetic
costs. Infinity must be handled explicitly, for example by a sufficiently
large sentinel in this nonnegative setting; it is not an ordinary
integer exponent in the next encoding.

\textbf{3. A correct algebraic encoding with the wrong size parameter.}
Set $b=n+1$ and encode the finite matrices as
\[
 U_{ik}=b^{A_{ik}},\qquad V_{kj}=b^{B_{kj}}.
\]
Then their ordinary product satisfies
\begin{equation}\label{30006803:base}
 (UV)_{ij}=\sum_{s=0}^{2W}c_{ij,s}b^s,
 \qquad c_{ij,s}=\#\{k:A_{ik}+B_{kj}=s\}\leq n<b.
\end{equation}
There are no carries. The least index of a nonzero base-$b$ digit is
exactly $(A\star B)_{ij}$. This proves correctness of the reduction; it
is not a probabilistic identity or an approximation.

However the input encodings $b^{A_{ik}}$ may require
$\Theta(W\log(n+1))$ bits, and outputs in \eqref{30006803:base} can require
that order of bits as well. The original entries required only
$O(\log(W+1))$ bits. If $W=2^B$, the representation has exponential
size in the original weight bit length $B$. Counting each multiplication
of those huge integers as one word operation suppresses precisely the
cost that must be justified. Fast ordinary matrix multiplication can
make this a useful bounded-magnitude method, but correctness of
\eqref{30006803:base} alone is not a truly subcubic algorithm with the
unrestricted intended weight convention.

Replacing the large integers by residues introduces another issue. The
first nonzero base digit is not generally determined by a few residues:
a nonzero coefficient contribution can vanish modulo the chosen modulus,
and its exponent can wrap around. A proof would need a deterministic
or controlled-error recovery procedure with total cost polynomial in
$\log W$ and the desired exponent in $n$. Such a recovery procedure is
not constructed here. The companion finite computation checks
\eqref{30006803:gadget} and \eqref{30006803:base} on exact small integer matrices;
it does not establish a complexity bound by timing them.

\textbf{4. Approximation-to-exact rounding loses the uniformity.}
Suppose an approximation algorithm guarantees
$D\leq\widetilde D\leq(1+\epsilon)D$ for an integer distance $D$.
Rounding to the nearest integer is certainly valid when
$\epsilon D<1/2$. In the gadget, $D\leq2H+2W=6W+2$,
so the sufficient choice is
\begin{equation}\label{30006803:rounding}
 \epsilon<\frac1{12W+4}.
\end{equation}
A guarantee for fixed $\epsilon>0$ cannot simply be invoked at this
input-dependent precision without rechecking the running time and bit
costs. Conversely, at relative error comparable to $1/W$, the valid
intervals for distances of order $W$ differing by one already begin to
overlap. The difficulty is not lack of an integer output convention;
it is the precision dependence. The approximation results in the
catalogue source retain their approximation parameters and do not assert
an exact algorithm after an arbitrary zero-error limit.
\eref{30006803}{1}.

A useful local stability observation does not remove the difficulty.
If the smallest of the candidate sums $A_{ik}+B_{kj}$ is separated
from every larger candidate by $\gamma>0$, and every candidate is
estimated with absolute error less than $\gamma/2$, then a minimizer
can be recovered by comparing estimates. This follows by subtracting
the two error bounds. But in the unrestricted integer problem one
may have ties, or a separation of only one next to entries of size
$W$. No fixed relative margin is provided. Perturbing weights to force
unique minimizers must also retain enough information to reconstruct
the original exact distances; uniqueness alone is not a speedup.

\paragraph{Stress test.}
The gadget explicitly excludes longer undirected detours by a strict
weight inequality. Using the naive three-layer graph without $H$ could
allow such detours and invalidate the reduction. The exponent encoding
uses base $n+1$, not base two; otherwise coefficient carries could change
the first nonzero digit. Both repairs are correctness repairs, not a
solution to the complexity obstacle.

The repeated-squaring implication is stated only for a product algorithm
that handles the enlarged weight range and unreachable pairs. It does
not import a bounded-weight theorem outside its hypotheses. Likewise,
$O(n^3/\log^c n)$ is not $O(n^{3-\eta})$ for any fixed $\eta>0$.
The two-vertex bit-output example is an objection to an unspecified
encoding, not a proof of a cubic lower bound in a standard RAM model.
No assertion of such a lower bound is made.

\needspace{15\baselineskip}
\paragraph{Outcome.}
\textbf{Outcome: Exploratory approach with unresolved gap.}

The attempt supplies an exact undirected positive-weight min-plus gadget,
an exact carry-free algebraic encoding, and explicit precision and
bit-length calculations showing where two proposed speedups stop.
These are reductions and diagnostics, not claimed new matrix-product
algorithms or lower bounds. The source's missing global machine convention
has been left visible throughout.

A resolution requires a genuinely faster exact min-plus mechanism with
the appropriate uniform weight/model guarantees, or an applicable lower
bound after the intended model has been fixed. Neither is obtained here.
The small exact checks verify only the displayed identities and gadgets.
% END RESEARCH ATTEMPT 232
\subsection{Sources}
\begin{itemize}\raggedright
\item[\textnormal{[S1]}]\hypertarget{source:30006803:1}{} Barna Saha and Christopher Ye, Faster Approximate All Pairs Shortest Paths, SODA 2024, Abstract.
\url{https://epubs.siam.org/doi/10.1137/1.9781611977912.170}.
{\small\textit{Catalogue formulation source.}}
% BEGIN ADDED REFERENCES 232
\item[\textnormal{[E1]}]\hypertarget{extra:30006803:1}{}
B. Saha and C. Ye, \emph{Faster Approximate All Pairs Shortest Paths},
SODA 2024; full manuscript arXiv:2309.13225v1, 23 September 2023,
Abstract, Introduction and the weighted-graph approximation sections.
\url{https://arxiv.org/html/2309.13225v1}.
Consulted 27 September 2026; the full manuscript supplements the
catalogue's proceedings abstract. Its approximation results and
historical matrix-multiplication exponent are not promoted to an exact
current algorithm.
% END ADDED REFERENCES 232
\end{itemize}

\end{document}
